%! Absolute Value Functions & Transformations — Unit 1, Lesson 1.3 — Guided Notes
% Gradual release (course-wide shape, 2026-09-01): the notes are the "I do / we do" block,
% 20 minutes, and the Individual Practice block that closes them is the "you do alone", 15 more.
% There is no group activity.
\documentclass[10pt]{article}
\usepackage{algebra2-article}
\usepackage{algebra2-boxes}
% Coordinate grid: {xmin}{xmax}{ymin}{ymax}
\newcommand{\lgrid}[4]{%
  \draw[step=1, linegray!40, very thin] (#1,#3) grid (#2,#4);
  \draw[->, semithick, charcoal] (#1,0)--(#2,0) node[below right, font=\scriptsize]{$x$};
  \draw[->, semithick, charcoal] (0,#3)--(0,#4) node[left, font=\scriptsize]{$y$};}

\begin{document}

\pageheader{Unit 1, Lesson 1.3}{Guided Notes}

\vspace{0.10in}

\begin{objectivebox}
\textbf{By the end of this lesson, I will be able to\dots}
\begin{itemize}\setlength{\itemsep}{2pt}
  \item read the \emph{absolute value parent} $f(x)=|x|$ as a V, and name its vertex, axis of
        symmetry, and minimum;
  \item use \emph{vertex form} $g(x)=a|x-h|+k$ to find the vertex $(h,k)$ and axis $x=h$ ---
        and say why the inside of the bars moves the graph \emph{backwards};
  \item decide from $a$ whether a V opens up (a \emph{minimum}) or down (a \emph{maximum}), and
        whether it is narrower or wider than the parent;
  \item read domain, range, and increasing/decreasing intervals off a V, and build its rule from
        a graph.
\end{itemize}
\end{objectivebox}

\vspace{0.05in}

\begin{vocabbox}
\small
Fill in each term as we build it together.
\par\vspace{2pt}   % \par is required: \termblanklong's \noindent is a no-op mid-paragraph
\termblanklong{Absolute value parent function}
\termblanklong{Vertex (turning point)}
\termblanklong{Axis of symmetry}
\termblanklong{Vertex form $a|x-h|+k$}
\termblanklong{Reflection}
\termblanklong{Vertical stretch / compression}
\end{vocabbox}

\begin{hookbox}
A trail drops toward a creek, crosses it on a bridge, and climbs the far side. Its height above
the creek, in \emph{hundreds of feet}, $x$ miles from the trailhead, is $H(x)=2|x-3|+1$.
\begin{center}
\begin{minipage}[c]{0.36\linewidth}\centering
\begin{tikzpicture}[scale=0.38]
  \lgrid{-0.5}{6.5}{-0.5}{8}
  \draw[very thick, forest] (0,7)--(3,1)--(6,7);
  \fill[charcoal] (3,1) circle (3.5pt);
  \node[charcoal, font=\scriptsize, below right] at (3,1){$(3,1)$};
\end{tikzpicture}
\end{minipage}\hfill
\begin{minipage}[c]{0.60\linewidth}
\begin{itemize}\setlength{\itemsep}{2pt}
  \item $H(0)=$ \blank{1.0cm} \quad $H(3)=$ \blank{1.0cm} \quad $H(6)=$ \blank{1.0cm}
  \item The trail is lowest at $x=$ \blank{0.9cm} miles, where the height is \blank{0.9cm}
        --- the bridge.
  \item Why can the height never drop below $1$? \blank{4.0cm}
\end{itemize}
\end{minipage}
\end{center}
Every absolute value function is this same \textbf{V} --- the parent $y=|x|$ after it has been
moved, stretched, or flipped. Today we read those moves straight off the rule.
\end{hookbox}

\boxguard[22]
\begin{notesbox}{1. The Parent $f(x)=|x|$ --- a V That Turns Once}
\begin{center}
\begin{minipage}[c]{0.30\linewidth}\centering
\renewcommand{\arraystretch}{1.15}
\begin{tabular}{|c|c|}
\hline
$x$ & $y=|x|$ \\ \hline
$-2$ & \blank{0.9cm} \\ \hline
$-1$ & \blank{0.9cm} \\ \hline
$0$ & \blank{0.9cm} \\ \hline
$1$ & \blank{0.9cm} \\ \hline
$2$ & \blank{0.9cm} \\ \hline
\end{tabular}
\end{minipage}\hfill
\begin{minipage}[c]{0.30\linewidth}\centering
\begin{tikzpicture}[scale=0.42]
  \lgrid{-4}{4}{-1}{5}
  \draw[very thick, forest] (-4,4)--(0,0)--(4,4);
  \fill[charcoal] (0,0) circle (3.5pt);
\end{tikzpicture}
\end{minipage}\hfill
\begin{minipage}[c]{0.34\linewidth}
\small
Because $|x|$ is a \emph{distance}, the outputs never go below \blank{0.7cm}, and $x$ and $-x$
give the \emph{same} output --- so the graph is a mirror image across the $y$-axis.
\end{minipage}
\end{center}
\textbf{Read the V:} vertex \blank{1.4cm}; \; axis of symmetry $x=$ \blank{0.8cm}; \; minimum
value \blank{0.8cm}; \; domain \blank{2.0cm}; \; range \blank{1.6cm}. It \emph{decreases} on
\blank{1.8cm} and \emph{increases} on \blank{1.8cm}. The V is really two lines pasted at the
vertex: $y=$ \blank{0.9cm} when $x\ge0$, and $y=$ \blank{0.9cm} when $x<0$.
\end{notesbox}

\boxguard[26]
\begin{notesbox}{2. Move It: Outside the Bars Is Vertical; Inside Is Horizontal --- and Backwards}
Adding \emph{outside} the bars moves the V up or down, as you expect. Adding \emph{inside} the
bars moves it left or right --- the \emph{opposite} way from the sign you see.
\begin{center}
\renewcommand{\arraystretch}{1.4}
\begin{tabularx}{\linewidth}{@{}l l X@{}}
\textbf{Rule} & \textbf{Moves the V} & \textbf{Examples} \\
\hline
$y=|x|+k$ & \blank{2.2cm} by $k$ & $y=|x|-3$ has vertex \blank{1.4cm} \\
$y=|x-h|$ & \blank{2.2cm} by $h$ & $y=|x-3|$ has vertex \blank{1.4cm}; \; $y=|x+2|$ has vertex \blank{1.4cm} \\
\end{tabularx}
\end{center}
Each solid V is the parent (dashed) after \emph{one} move. Write each vertex on the line below.
\begin{center}
\begin{tikzpicture}[scale=0.38]
  \begin{scope}
    \lgrid{-2}{6}{-1}{5}
    \draw[very thick, forest!40, dashed] (-1,1)--(0,0)--(4,4);
    \draw[very thick, forest] (-1,4)--(3,0)--(6,3);
    \fill[charcoal] (3,0) circle (3.5pt);
    \node[charcoal, font=\scriptsize] at (2,-2.2){$y=|x-3|$};
  \end{scope}
  \begin{scope}[xshift=11cm]
    \lgrid{-4}{4}{-4}{3}
    \draw[very thick, forest!40, dashed] (-3,3)--(0,0)--(3,3);
    \draw[very thick, forest] (-3,0)--(0,-3)--(3,0);
    \fill[charcoal] (0,-3) circle (3.5pt);
    \node[charcoal, font=\scriptsize] at (0,-5.2){$y=|x|-3$};
  \end{scope}
\end{tikzpicture}
\end{center}
Left vertex: \blank{1.4cm} \qquad Right vertex: \blank{1.4cm}
\begin{tcolorbox}[colback=goldbg, colframe=goldacc, boxrule=0.9pt, arc=2mm,
    left=2.5mm, right=2.5mm, top=1.5mm, bottom=1.5mm]
\small \textbf{Find the vertex by setting the inside to zero.} $|x-3|$ is smallest when
$x-3=0$, that is at $x=$ \blank{0.8cm} --- so $|x-3|$ moves the V \emph{right}, not left. For the
trail $H(x)=2|x-3|+1$: inside $=0$ at $x=$ \blank{0.8cm}, and the $+1$ outside lifts it, so the
vertex is \blank{1.4cm} --- the bridge from the hook.
\end{tcolorbox}
\end{notesbox}

\boxguard[26]
\begin{notesbox}{3. The Number $a$: Up or Down, Narrow or Wide}
The trail is a \emph{valley}. Change the sign of $a$ and the same vertex becomes a \emph{ridge}.
\begin{center}
\renewcommand{\arraystretch}{1.35}
\begin{tabularx}{\linewidth}{@{}X|X@{}}
\textbf{Valley: $H(x)=2|x-3|+1$} & \textbf{Ridge: $R(x)=-2|x-3|+7$} \\
\hline
\centering\arraybackslash
\begin{tikzpicture}[scale=0.30, baseline=(current bounding box.center)]
  \lgrid{-0.5}{6.5}{-0.5}{8}
  \draw[very thick, forest] (0,7)--(3,1)--(6,7);
  \fill[charcoal] (3,1) circle (4pt);
\end{tikzpicture}
&
\centering\arraybackslash
\begin{tikzpicture}[scale=0.30, baseline=(current bounding box.center)]
  \lgrid{-0.5}{6.5}{-0.5}{8}
  \draw[very thick, redacc] (0,1)--(3,7)--(6,1);
  \fill[charcoal] (3,7) circle (4pt);
\end{tikzpicture}
\\
$a=2>0$: opens \blank{1.2cm} & $a=-2<0$: opens \blank{1.2cm} --- a \emph{reflection} \\
Vertex \blank{1.2cm}; the value $1$ is a \blank{1.6cm} & Vertex \blank{1.2cm}; the value $7$ is a \blank{1.6cm} \\
Range: $y$ \blank{0.8cm} $1$ & Range: $y$ \blank{0.8cm} $7$ \\
Decreasing on \blank{1.6cm}, increasing on \blank{1.6cm} & Increasing on \blank{1.6cm}, decreasing on \blank{1.6cm} \\
\end{tabularx}
\end{center}
\vspace{2pt}
\begin{tcolorbox}[colback=goldbg, colframe=goldacc, boxrule=0.9pt, arc=2mm,
    left=2.5mm, right=2.5mm, top=1.5mm, bottom=1.5mm]
\small \textbf{Careful --- a negative $a$ does not just flip the picture; it flips the
vertex's job.} A classmate says ``the ridge still has a \emph{minimum} at its vertex, $7$.''
Test $x=5$: $R(5)=-2|5-3|+7=$ \blank{0.9cm}. Is that above or below $7$? \blank{1.2cm} \; So $7$
is the \blank{1.4cm} (highest / lowest) output, a \textbf{maximum}, and every other output sits
\emph{below} it. Same vertex, opposite range.
\end{tcolorbox}
\vspace{2pt}
\textbf{The size of $a$ sets the width.} $a$ is the \emph{slope of the right-hand ray}: from the
valley's vertex $(3,1)$, one step right rises to $(4,3)$ --- slope \blank{0.8cm}. So $|a|>1$ makes
the V \blank{1.8cm} than the parent, and $0<|a|<1$ makes it \blank{1.8cm}: $y=2|x|$ is narrower;
$y=\tfrac12|x|$ is wider.
\end{notesbox}

\boxguard[24]
\begin{notesbox}{4. Vertex Form, All at Once --- and Back to Lesson 1.2}
\[ g(x)=a|x-h|+k: \qquad \text{vertex } (h,k), \qquad \text{axis } x=h \]
\vspace{-6pt}
\begin{center}\small $a>0$: opens up, $k$ a \textbf{minimum} \qquad $a<0$: opens down, $k$ a \textbf{maximum}\end{center}
\textbf{Worked: read every feature of $g(x)=-2|x-1|+3$ from the rule.}
\begin{center}
\renewcommand{\arraystretch}{1.4}
\begin{tabularx}{\linewidth}{@{}X X X X X@{}}
Vertex \blank{1.2cm} & Axis $x=$ \blank{0.7cm} & Opens \blank{1.0cm} & \blank{0.9cm} value $3$ & Range \blank{1.2cm} \\
\end{tabularx}
\end{center}
Compared to $y=|x|$: shifted \blank{1.2cm} $1$ and \blank{1.2cm} $3$, reflected, and
\blank{1.5cm} (narrower / wider). Evaluate $g(0)$ straight from the rule, one step per line:
\begin{work}
  g(0) &= -2\,|0-1| + 3 \\
       &= -2(1) + 3 \\
       &= 1
\end{work}
\textbf{Back to Lesson 1.2.} Where is the trail exactly $500$ ft up? That is $H(x)=5$:
\begin{work}
  2|x-3| + 1 &= 5 \\
  |x-3| &= 2 \\
  x - 3 &= 2 \quad\text{or}\quad x - 3 = -2 \\
  x &= 5 \quad\text{or}\quad x = 1
\end{work}
The two solutions you solved for last lesson are the two places the V crosses the horizontal
line $y=5$ --- one on each side of the axis $x=$ \blank{0.8cm}.
\end{notesbox}

\boxguard[20]
\begin{practicebox}
For $g(x)=\tfrac12\,|x+4|-2$, read every feature; then build a rule from a graph.
\begin{enumerate}[leftmargin=1.7em, itemsep=3pt]
  \item Vertex \blank{1.4cm}; \; axis $x=$ \blank{0.9cm}; \; opens \blank{1.2cm}, so $-2$ is a
        \blank{1.4cm}; \; range \blank{1.6cm}; \; \blank{1.4cm} than the parent.
  \item Evaluate $g(0)$, one step per line.
        \begin{work}
          g(0) &= \tfrac12\,|0+4| - 2 \\
               &= \tfrac12(4) - 2 \\
               &= 0
        \end{work}
  \item \textbf{Build the rule from the graph.} \\[2pt]
        \begin{minipage}[c]{0.62\linewidth}
        Vertex \blank{1.4cm}. From the vertex, the right-hand ray reaches $(1,0)$: it falls
        \blank{0.7cm} over a run of \blank{0.7cm}, so $a=$ \blank{0.9cm}. \\[3pt]
        Rule: $y=$ \blank{3.2cm}
        \end{minipage}\hfill
        \begin{minipage}[c]{0.34\linewidth}\centering
        \begin{tikzpicture}[scale=0.36]
          \lgrid{-4}{2}{-1}{5}
          \draw[very thick, forest] (-3,0)--(-1,4)--(1,0);
          \fill[charcoal] (-1,4) circle (4pt);
          \fill[charcoal] (1,0) circle (3pt);
        \end{tikzpicture}
        \end{minipage}
\end{enumerate}
\end{practicebox}

\boxguard[22]
\begin{scenariobox}[Individual Practice --- On Your Own]{navy}
Three problems, quietly and on your own. Find the vertex by setting the inside to zero, then let
the sign of $a$ tell you up or down.
\begin{enumerate}[leftmargin=1.9em, itemsep=9pt]
  \item \textbf{From the rule.} $p(x)=-3|x-2|+5$. \\[2pt]
    Vertex \blank{1.4cm}; \; axis $x=$ \blank{0.9cm}; \; opens \blank{1.2cm}; \; the value $5$ is a
    \blank{1.5cm}; \; range \blank{1.6cm}; \; \blank{1.5cm} than the parent. \\[2pt]
    Evaluate $p(4)$:
    \begin{work}
      p(4) &= -3\,|4-2| + 5 \\
           &= -3(2) + 5 \\
           &= -1
    \end{work}

  \item \textbf{From the graph.} \\[2pt]
    \begin{minipage}[c]{0.62\linewidth}
    Vertex \blank{1.4cm}; \; axis $x=$ \blank{0.9cm}. From the vertex to the point $(3,2)$ the
    right-hand ray rises \blank{0.7cm} over a run of \blank{0.7cm}, so $a=$ \blank{0.9cm}. \\[3pt]
    Rule: $y=$ \blank{3.0cm} \\[3pt]
    The V crosses the dashed line $y=2$ at $x=$ \blank{0.8cm} and $x=$ \blank{0.8cm}.
    \end{minipage}\hfill
    \begin{minipage}[c]{0.34\linewidth}\centering
    \begin{tikzpicture}[scale=0.36]
      \lgrid{-3}{5}{-3}{4}
      \draw[dashed, redacc] (-3,2)--(5,2);
      \draw[very thick, forest] (-2,4)--(1,-2)--(4,4);
      \fill[charcoal] (1,-2) circle (4pt);
      \fill[charcoal] (3,2) circle (3pt);
    \end{tikzpicture}
    \end{minipage} \\[2pt]
    Check those two crossings by solving your rule $=2$:
    \begin{work}
      2|x-1| - 2 &= 2 \\
      |x-1| &= 2 \\
      x - 1 &= 2 \quad\text{or}\quad x - 1 = -2 \\
      x &= 3 \quad\text{or}\quad x = -1
    \end{work}

  \item \textbf{Same vertex, opposite job.} Compare $y=|x+2|-3$ and $y=-|x+2|-3$. \\[2pt]
    Both have vertex \blank{1.4cm} and axis $x=$ \blank{0.9cm}. The first opens \blank{1.0cm}
    with a \blank{1.4cm} of $-3$ and range \blank{1.5cm}; the second opens \blank{1.0cm} with a
    \blank{1.4cm} of $-3$ and range \blank{1.5cm}. \\[2pt]
    A classmate says ``both have a minimum of $-3$.'' What did they miss?
    \par\writelines{2}
\end{enumerate}
\end{scenariobox}

\end{document}
